\section{Toroidal topology and validated Casimir backreaction}

\subsection{Periodic rectangular lattice}

For a massless scalar with periodic lengths $(L_1,L_2,L_3)$, the renormalized
energy density used in CBR-001 is the rectangular-periodic Epstein-lattice
result developed in Ref.~\cite{edery2006}:
\begin{equation}
 \rho_{\rm Cas}=-\frac{\hbar c}{2\pi^2}
 \sum_{\mathbf n\in\mathbb Z^3\setminus\{0\}}
 \frac{1}{\left[(n_1L_1)^2+(n_2L_2)^2+(n_3L_3)^2\right]^2}.
 \label{eq:casimir-lattice}
\end{equation}
The directional pressures follow thermodynamically,
\begin{equation}
 p_i=-\frac{1}{L_jL_k}\frac{\partial E_{\rm Cas}}{\partial L_i},
 \qquad E_{\rm Cas}=L_1L_2L_3\rho_{\rm Cas}.
 \label{eq:casimir-pressure}
\end{equation}
The Stage-1 and Stage-2 calculations validate cutoff convergence, cubic
isotropy, pressure derivatives, the massless trace relation, dimensional
scaling and the sign change of $p_t-p_p$ across the cubic shape.  Thus compact
shape-dependent anisotropic stress is a concrete mechanism.  It is not a
universal cycle-counting coefficient.  A recent direct semiclassical
calculation likewise derives a finite topology-dependent stress tensor on
$T^3$ and its anisotropic de Sitter backreaction~\cite{negro2026}.

\subsection{Biaxial backreaction}

For
\begin{equation}
 \dd s^2=-\dd t^2+a_p^2\dd x^2+a_t^2(\dd y^2+\dd z^2),
\end{equation}
define
\begin{equation}
 H=\frac{H_p+2H_t}{3},
 \quad
 \delta=H_t-H_p,
 \quad
 r=\frac{a_t}{a_p}.
\end{equation}
Then
\begin{equation}
 \dot\delta+3H\delta=\kappa(p_t-p_p),
 \qquad
 \dot r=r\delta,
 \label{eq:shear-system}
\end{equation}
with the Hamiltonian constraint
\begin{equation}
 3H^2-\frac{\delta^2}{3}=\kappa\rho_{\rm total}.
 \label{eq:hamiltonian-biaxial}
\end{equation}
These are the biaxial specialization of the standard Bianchi-I system; see,
for example, Ref.~\cite{ellismaccallum1969}.
The Stage-3A solver reproduces the perturbative solution and preserves the
constraint and Casimir continuity equation to substantially better than the
specified tolerances.

\subsection{The Stage-3B ratio test}

Let
\begin{equation}
 x=\frac{\delta}{H},
 \qquad
 q=\frac{H_t}{H_p}=\frac{1+x/3}{1-2x/3}.
 \label{eq:q-x}
\end{equation}
The historical target $q=13/12$ corresponds to
\begin{equation}
 x_\star=\frac{3}{38}.
\end{equation}
The small-source fixed-background solution is
\begin{equation}
 x(N)=S\left(e^{-3N}-e^{-4N}\right),
\end{equation}
with maximum $(27/256)S$ at $N=\ln(4/3)$.  Merely reaching $x_\star$ therefore
requires
\begin{equation}
 S\ge\frac{128}{171}\simeq0.748538,
\end{equation}
already outside the strict small-perturbation regime.

The nonlinear root search finds target-reaching amplitudes for five of the
tested initial shapes.  Every case is classified as a transient crossing,
with only approximately $0.27$--$0.30$ e-folds within one percent of the
target.  Four thresholds are nonperturbative and one is marginal by the stated
Casimir-fraction diagnostic.  Every valid candidate returns to $q=1$ by
$N=20$; no quasi-plateau or attractor is found.

There is also a kinematic warning.  Since
\begin{equation}
 \frac{\dd\ln r}{\dd N}=x,
\end{equation}
a constant $q\ne1$ implies continuous evolution of the shape ratio.  It cannot
be a fixed torus shape.  A persistent anisotropic state would have to be a
driven solution, not a static geometric ratio.

For a desired trajectory $x(N)$, the required anisotropic stress is
\begin{equation}
 \Pi_{\rm req}=\frac{H^2}{\kappa}
 \left[\frac{\dd x}{\dd N}
 +x\left(3+\frac{\dd\ln H}{\dd N}\right)\right].
 \label{eq:required-pi}
\end{equation}
In de Sitter with constant $x=3/38$ this becomes
\begin{equation}
 \Pi_{\rm req}=\frac{9H^2}{38\kappa}.
\end{equation}
For the CBR-001 physical lengths
$L_p=L_\star a r^{-2/3}$ and $L_t=L_\star a r^{1/3}$, the anisotropic Casimir
source is more precisely
\begin{equation}
 \Pi_{\rm Cas}(a,r)=\frac{\hbar c}{L_\star^4}a^{-4}
 r^{8/3}\left[\widehat p_t(r)-\widehat p_p(r)\right]
 \equiv\frac{\hbar c}{L_\star^4}a^{-4}F_\Pi(r).
 \label{eq:casimir-pi-scaling}
\end{equation}
Thus the explicit scale dependence is radiation-like, while the shape
function evolves whenever $r$ evolves.  In every validated free-field
trajectory the source nevertheless becomes negligible and cannot maintain
Eq.~\eqref{eq:required-pi}.  A future CBR-002 test must derive any additional
shape, driven or memory stress without inserting the target ratio into a
coefficient.
